% =========================================================================
% Do Errors Compound? Measuring Error Propagation in Multi-Agent LLM Pipelines
% SISTER 2026 --- AI/ML Track (Dubai CS Society + Delta Rising Foundation)
%
% Self-contained: compiles on Overleaf as-is (pdfLaTeX). When moving to the
% official NeurIPS template later, only the preamble changes; the body is
% written to NeurIPS register already.
%
% Convention: \pending{...} marks every number or claim that must come from
% our own runs. Nothing pending may survive into a submitted version.
% =========================================================================
\documentclass[11pt]{article}
\usepackage[margin=1in]{geometry}
\usepackage{amsmath,amssymb}
\usepackage{booktabs}
\usepackage{graphicx}
\usepackage{microtype} % NOTE: if a local build errors on font expansion, use [expansion=false]
\usepackage[hidelinks]{hyperref}
\usepackage[numbers,sort&compress]{natbib}
\usepackage{xcolor}
\newcommand{\pending}[1]{\textcolor{purple}{[#1]}}

\title{\bf Do Errors Compound? Measuring Error Propagation\\in Multi-Agent LLM Pipelines}
\author{Tilak Parajuli\\
\small SISTER 2026, AI/ML Track --- Dubai Computer Science Society \& Delta Rising Foundation\\
\small Kathmandu, Nepal
\and \pending{Teammate Name(s)}\\
\small \pending{affiliation}}
\date{August 2026}

\begin{document}
\maketitle

\begin{abstract}
Pipelines of specialized LLM agents---one agent plans, another executes, a
third verifies---are now a default design for applied language-model systems.
The design carries an implicit promise: later stages should absorb the
mistakes of earlier ones. Whether they actually do is unresolved. Failure
taxonomies built from deployed systems describe \emph{what} goes wrong but,
being observational, cannot say how often an early error survives to the
output or what a verification stage causally contributes. We address this
with controlled fault injection, adapted from software reliability
engineering: we corrupt a single intermediate output with a known, scripted
perturbation and trace its fate through the rest of the pipeline. Across
\pending{200} GSM8K problems \pending{and 150 HotpotQA questions}, we compare
a single-agent baseline, a clean three-stage pipeline, early- and late-stage
injection, and a verifier-removed ablation, analyzing outcomes with paired
statistics throughout. Injected faults reached the final answer in
\pending{XX\%} of cases; the verifier flagged \pending{YY\%} of faults but
repaired only \pending{ZZ\%}, and faults planted at the planning stage
survived \pending{significantly more often} than those planted during
execution. We release code, prompts, and full traces.
\end{abstract}

\section{Introduction}

The unit of deployment for language models is shifting from a single model
call to a pipeline of agents. A planning agent decomposes the task, a solving
agent executes the plan, and a verification agent checks the result before it
leaves the system; frameworks such as AutoGen~\citep{wu2023autogen} and
MetaGPT~\citep{hong2023metagpt} have made this pattern easy to build and
correspondingly common. The pattern borrows its justification from human
organizations: specialization should raise quality, and review should catch
what specialization misses.

The evidence for that justification is mixed. On one side, multi-agent debate
has been reported to improve factuality and reasoning~\citep{du2023debate},
and simple ensembling over agents continues to help as the number of agents
grows~\citep{li2024moreagents}. On the other, a growing body of work finds
that language models are poor judges of their own reasoning: without external
feedback, self-correction often leaves accuracy unchanged or makes it
worse~\citep{huang2023selfcorrect}. Between these positions sits the most
detailed account of multi-agent behavior to date, the MAST
taxonomy~\citep{cemri2025mast}, which catalogues fourteen failure modes
observed across seven deployed frameworks. MAST establishes that inter-agent
failures are frequent and varied. But it is an observational study: when a
deployed system fails, the origin of the failure must be reconstructed after
the fact, and how often a given early-stage error would have survived to the
output is not identifiable from logs alone.

What is missing is an interventional measurement. If we could plant a known
error at a chosen stage and watch what the rest of the pipeline does with it,
questions that the observational work leaves open would become directly
answerable: How often does an early mistake reach the user? Does it matter
where in the pipeline the mistake occurs? How much of the damage does a
verifier actually prevent, as opposed to merely being present while the
pipeline succeeds for other reasons?

We make that measurement. Borrowing fault injection from software
reliability engineering, we corrupt exactly one intermediate artifact of a
Decomposer--Solver--Verifier pipeline---swapping a quantity or inverting an
arithmetic operation---and let the remaining stages run untouched. Because
the fault is planted by us, its origin is known by construction, and every
downstream outcome can be attributed. On fixed problem samples from
GSM8K~\citep{cobbe2021gsm8k} \pending{and HotpotQA~\citep{yang2018hotpotqa}},
we evaluate five conditions: a single-agent baseline, the clean pipeline,
injection at the planning stage, injection at the solving stage, and the
injected pipeline with its verifier removed. All comparisons are paired, and
we report exact tests and bootstrap intervals rather than point estimates
alone.

\textbf{Contributions.} (i)~A controlled fault-injection methodology for
measuring error propagation in LLM agent pipelines, with scripted,
reproducible perturbation operators. (ii)~Quantitative propagation, catch,
and correction rates, including the effect of injection position, under
paired statistical tests. (iii)~A verifier ablation that isolates the causal
contribution of the verification stage. \pending{(iv)~A 1.5B-parameter
verifier fine-tuned with QLoRA~\citep{dettmers2023qlora} on our own traces,
compared against an LLM judge on catch rate and cost.} (v)~Code, prompts,
configuration, and complete traces, sufficient to regenerate every number in
this paper from one script.

\section{Related work}

\textbf{Multi-agent frameworks.} AutoGen~\citep{wu2023autogen},
MetaGPT~\citep{hong2023metagpt}, and CAMEL~\citep{li2023camel} established
conversational and assembly-line patterns for composing LLM agents; the
Decomposer--Solver--Verifier pipeline we study is the smallest configuration
that exhibits the planning, execution, and review roles these frameworks
popularized. Our contribution is not a new framework but a measurement of a
failure property this family of systems shares.

\textbf{Does adding agents help?} Multi-agent debate improves factuality on
several benchmarks~\citep{du2023debate}, and majority voting over sampled
agents scales favorably~\citep{li2024moreagents}; self-consistency achieves
related gains within a single model~\citep{wang2022selfconsistency}. These
results measure aggregate accuracy. They do not distinguish a pipeline that
prevents errors from one that merely commits different ones, which is
precisely the distinction fault injection makes observable.

\textbf{Self-correction and verification.} Reflexion~\citep{shinn2023reflexion}
and Self-Refine~\citep{madaan2023selfrefine} show that models can revise
their outputs given feedback, while \citet{huang2023selfcorrect} demonstrate
that without external signals such revision often fails on reasoning tasks.
LLM-as-judge evaluation is standard but carries known
biases~\citep{zheng2023judge}. Our verifier ablation speaks to this debate
directly: it estimates, under identical inputs, how much a review stage
changes the fate of a known error.

\textbf{Failure analysis.} MAST~\citep{cemri2025mast} provides the
vocabulary---specification failures, inter-agent misalignment, verification
failures---and we adopt it. The difference is method: MAST classifies
failures that occurred in the wild; we induce failures under control. The two
are complementary in the way epidemiology and randomized trials are.

\section{Method}

\subsection{Pipeline under study}
Let $x$ be a problem. The pipeline computes a plan $d = D(x)$, a solution
$s = S(x, d)$, and a verdict $v = V(x, s)$; the final answer $\hat{y}$ is
extracted from $s$ if the verifier accepts, or from the verifier's correction
if it rejects. Each stage is a separate LLM call with its own role prompt
(Appendix~A); the only information passed between stages is the text of the
preceding artifact, mirroring common framework defaults. A single-agent
baseline computes $\hat{y}$ from one call on $x$ alone.

\begin{figure}[t]
\centering
% \includegraphics[width=.9\linewidth]{figures/fig1_pipeline.pdf}
\caption{The Decomposer--Solver--Verifier pipeline with the two injection
points marked. \pending{Figure to be rendered.}}
\label{fig:pipeline}
\end{figure}

\subsection{Controlled fault injection}
An injector $\phi$ maps an artifact to a corrupted artifact and a record of
the change. We use two operators chosen to mimic the arithmetic slips that
occur naturally in reasoning traces: \emph{number swap}, which replaces one
quantity with a plausible wrong value (scaling by 0.4--1.7$\times$, or
$\pm1$--$3$ for small integers), and \emph{operation flip}, which inverts one
arithmetic operation word (\emph{add}$\leftrightarrow$\emph{subtract},
\emph{multiply}$\leftrightarrow$\emph{divide}, and analogous pairs). List and
step labels are excluded from number swap: corrupting ``Step 1'' changes no
quantity, and admitting such perturbations would dilute the injected-fault
population with harmless cases and bias propagation rates downward. Every
injection logs the operator, the original and replacement spans, and the
character position, so each downstream error is attributable to a specific,
inspectable change.

In condition C1 the injector is applied to the plan $d$ before the solver
sees it; in C2, to the solution $s$ before the verifier sees it. Exactly one
artifact is corrupted per run.

\subsection{Conditions and design}
\begin{table}[h]
\centering\small
\begin{tabular}{lccc}
\toprule
Condition & Pipeline & Injection & Verifier \\
\midrule
A  & --         & --          & --          \\
B  & \checkmark & --          & \checkmark  \\
C1 & \checkmark & Decomposer  & \checkmark  \\
C2 & \checkmark & Solver      & \checkmark  \\
D1 & \checkmark & Decomposer  & --          \\
D2 & \checkmark & Solver      & --          \\
\bottomrule
\end{tabular}
\caption{Experimental conditions. Every condition runs on the same fixed
problem sample (seeded), so all comparisons are paired at the problem level.}
\label{tab:conditions}
\end{table}

\subsection{Outcome measures}
For a problem the clean pipeline solves correctly, an injected fault has one
of four fates: \emph{caught and corrected} (verifier rejects; final answer
correct), \emph{caught but not corrected} (verifier rejects; final answer
wrong), \emph{silently propagated} (verifier accepts; final answer wrong), or
\emph{absorbed} (verifier accepts; final answer nevertheless correct, e.g.,
the solver ignored the corrupted span). We report: accuracy per condition;
\emph{propagation rate}, the fraction of valid injections whose final answer
is wrong, restricted to problems condition B solved correctly so that new
errors are attributable to the injection; \emph{catch rate} and
\emph{correction rate} of the verifier; and cost per solved problem in
tokens.

\subsection{Statistical analysis}
All condition contrasts use McNemar's exact test on paired
correct/incorrect outcomes~\citep{dietterich1998approximate}, the appropriate
test when two systems answer the same items. Accuracies carry 95\%
percentile-bootstrap intervals (10{,}000 resamples); rates carry Wilson
score intervals. The effect of injection position and task type on fault
survival is estimated by logistic regression, reported as odds ratios.
Because several contrasts are tested, $p$-values are Holm--Bonferroni
corrected~\citep{dror2018hitchhiker}. Main runs use temperature~0 on a fixed
seeded sample; a robustness replication uses temperature~0.7 over three
seeds, reported as mean and interval.

\section{Experimental setup}
Model: \pending{exact model string}, accessed \pending{dates}, temperature 0
(main), 0.7 (robustness). Datasets: GSM8K test split, $n=\pending{200}$
problems sampled with seed 42 \pending{; HotpotQA distractor setting,
$n=150$}. Injection seed 7; operators as in \S3.2. Prompts are fixed across
conditions and listed verbatim in Appendix~A. Total API cost:
\pending{\$XX}. All configuration is a single YAML file in the released
repository; \texttt{python -m src.stats} regenerates every number below from
the released traces.

\section{Results}
\pending{This section is written after the runs. Planned presentation:}

\textbf{Does the pipeline beat one agent?} Figure~\ref{fig:acc} compares
conditions A and B \pending{(bar chart, 95\% CIs; McNemar $p$)}.

\textbf{What happens to an injected fault?} Figure~\ref{fig:fate} gives the
error-fate breakdown for C1 and C2 \pending{(stacked bars or Sankey:
caught-and-corrected / caught-not-corrected / silently-propagated /
absorbed)}.

\textbf{Does position matter?} \pending{Propagation rate by injection stage
with odds ratio from the logistic model (H2).}

\textbf{What does the verifier buy?} \pending{C1 vs.\ D1 and C2 vs.\ D2:
propagation with and without the verifier; the difference is the verifier's
causal contribution.}

\begin{figure}[t]\centering
\caption{Accuracy by condition with 95\% bootstrap intervals.
\pending{pending}}\label{fig:acc}
\end{figure}
\begin{figure}[t]\centering
\caption{Fate of injected faults by injection stage. \pending{pending}}
\label{fig:fate}
\end{figure}

\section{Analysis}
\pending{Three short case studies drawn from traces: one silent propagation
(quote the corrupted span and the solver's uncritical use of it), one
catch-and-correct, one absorption. Each at most a paragraph; the point is
mechanism, not anecdote volume.}

\section{Limitations}
Our perturbations are synthetic by design; this is what makes origin
attribution exact, but naturally arising errors are more varied, and our
operators cover arithmetic slips rather than, e.g., factual hallucinations.
Results are reported for \pending{one model family} and \pending{one to two}
task distributions with objectively checkable answers; propagation dynamics
on open-ended tasks may differ. The verifier prompt is representative rather
than exhaustively tuned, so our estimate of verification benefit should be
read as typical-configuration, not best-achievable. The study was designed
and executed within a four-week program.

\section{Conclusion}
\pending{One paragraph, written last: what the numbers established, one
sentence on what changes for practitioners (where to spend verification
budget), one sentence on next steps.}

\paragraph{Broader impact.} Agent pipelines increasingly mediate consequential
work. Measuring where errors survive---and how much verification actually
helps---informs where guardrails belong; we do not foresee direct misuse
enabled by these measurements.

\paragraph{Reproducibility.} Code, prompts, configuration, and full traces:
\pending{repository URL}.

\bibliographystyle{plainnat}
\bibliography{references}

\appendix
\section{Role prompts}
\textbf{Decomposer.} \emph{You are a planning agent. Break the problem into a
short numbered plan of 2--5 concrete steps with any intermediate quantities
needed. Do NOT solve it. Output only the plan.}

\textbf{Solver.} \emph{You are a solving agent. Execute the given plan step by
step on the problem. Show brief working, then end with the line: FINAL
ANSWER: $\langle$answer$\rangle$.}

\textbf{Verifier.} \emph{You are a verification agent. Check the proposed
solution against the original problem. If correct, reply `VERDICT: ACCEPT'.
If you find an error, reply `VERDICT: REJECT', explain the error in one
sentence, and give a corrected line: FINAL ANSWER: $\langle$answer$\rangle$.}

\textbf{Single agent.} \emph{Solve the problem step by step. End with the
line: FINAL ANSWER: $\langle$answer$\rangle$.}

\section{Injection examples}
Number swap: ``Step 1: add \underline{45} and 12.'' $\rightarrow$ ``Step 1:
add \underline{76} and 12.'' (step labels untouched). ``Take the
\underline{8} apples'' $\rightarrow$ ``Take the \underline{3} apples.''
``Compute \underline{250} $\times$ 0.2'' $\rightarrow$ ``Compute
\underline{150} $\times$ 0.2.'' Operation flip: ``\underline{multiply} the
total by 3'' $\rightarrow$ ``\underline{divide} the total by 3.''

\end{document}
